\chapter{Phase 2: SIP test with a softphone}

\section{Goal}

Validate the Asterisk PBX before touching the HT801. To do so, a \emph{softphone} (Linphone) is registered as extension 102, the number 9000 is dialed, and we confirm that the DTMF tones sent from the app reach Asterisk as AMI events.

This step is important because it lets us verify the software side of the system without yet depending on the analog hardware, simplifying troubleshooting.

\section{Step 1: Install Linphone}

Linphone is a free, cross-platform SIP client:

\begin{itemize}
  \item \textbf{iOS:} App Store, search for \texttt{Linphone}.
  \item \textbf{Android:} Google Play, search for \texttt{Linphone}.
  \item \textbf{macOS:} \texttt{linphone.org/technical-corner/linphone}.
\end{itemize}

Install the app on any of the available devices that share the same local network as the Pi.

\section{Step 2: Create a SIP account in Linphone}

Open Linphone and enter the wizard: \textsf{Assistant} $\to$ \textsf{Use a SIP account} (or \textsf{Add Account}, depending on the version).

Fill in:

\begin{table}[h]
\centering
\begin{tabular}{ll}
\toprule
\textbf{Field} & \textbf{Value} \\
\midrule
Username & 102 \\
Password & the secret defined for extension 102 in \texttt{sip.conf} \\
Domain & 192.168.1.10 (the Pi's IP) \\
Transport & UDP \\
Display Name & Test Hugo \\
\bottomrule
\end{tabular}
\end{table}

Confirm and verify that the status indicator turns green with the legend \texttt{Registered} or \texttt{Connected}.

\section{Step 3: Verify the registration from Asterisk}

On the Pi:

\begin{shellcode}
sudo asterisk -rvvv
\end{shellcode}

Inside the CLI:

\begin{plaincode}
hugo-pbx*CLI> sip show peers
Name/username             Host                Status
101/101                   (Unspecified)       UNKNOWN
102/102                   192.168.1.20:xxxx   OK (3 ms)
\end{plaincode}

Extension 102 must appear with the IP of the device running Linphone and status \texttt{OK}.

\section{Step 4: Test call to 9000}

From Linphone, dial \texttt{9000} and press the call button.

Two things should be observed simultaneously:

\begin{enumerate}
  \item In the device's earpiece: a short \emph{beep}, produced by the \texttt{Playback(beep)} step in the dialplan.
  \item In the Asterisk console: the dialplan execution trace.
\end{enumerate}

\begin{plaincode}
-- Executing [9000@hugo-game:1] NoOp("PJSIP/102-...", "Hugo Phone — ...")
-- Executing [9000@hugo-game:2] Answer("PJSIP/102-...", "")
-- Executing [9000@hugo-game:3] Wait("PJSIP/102-...", "1")
-- Executing [9000@hugo-game:4] Playback("PJSIP/102-...", "beep")
-- <PJSIP/102-...> Playing 'beep.gsm' (language 'en')
-- Executing [9000@hugo-game:5] Read("PJSIP/102-...", "TECLA,,1,,,3600")
\end{plaincode}

\section{Step 5: Test DTMF}

While the call is still active, in Linphone open the tone keypad and press the digit \textbf{2}.

In the Asterisk console you should see:

\begin{plaincode}
== Manager 'hugo-bridge' logged on from 127.0.0.1
== UserEvent: HugoKey
   Channel: PJSIP/102-00000001
   Caller: 102
   Key: 2
-- Executing [9000@hugo-game:6] GotoIf(...)
-- Executing [9000@hugo-game:5] Read("PJSIP/102-...", "TECLA,,1,,,3600")
\end{plaincode}

Repeat with \textbf{4}, \textbf{6}, \textbf{8} and \textbf{5}. Each must trigger its own \texttt{UserEvent}.

If the events appear, \textbf{Phase 2 is complete}: the PBX captures DTMF correctly and the bridge will be able to read them.

\section{Troubleshooting}

\subsection{Linphone does not register}

\begin{itemize}
  \item Check that the Pi is reachable from the device: \texttt{ping 192.168.1.10}.
  \item Confirm that UDP port 5060 isn't blocked by a firewall on the Pi (Raspberry Pi OS Lite ships with no firewall enabled by default).
  \item Look at Asterisk logs for failed attempts:
  \begin{shellcode}
sudo asterisk -rvvv
hugo-pbx*CLI> sip set debug on
  \end{shellcode}
\end{itemize}

\subsection{The 9000 call drops immediately}

\begin{itemize}
  \item Check that extension 102 is assigned to context \texttt{hugo-game} in \texttt{sip.conf}.
  \item Run \texttt{dialplan show hugo-game} in the CLI and confirm that extension 9000 is present.
\end{itemize}

\subsection{It connects but the beep is not heard}

\begin{itemize}
  \item Typical NAT or incompatible-codec problem. Verify that \texttt{allow=ulaw} and \texttt{allow=alaw} are in \texttt{sip.conf} and that Linphone also supports these codecs in its advanced configuration.
  \item Check the device's volume.
\end{itemize}

\subsection{No UserEvent appears when sending DTMF}

\begin{itemize}
  \item Linphone may send DTMF as \emph{inband} tones instead of RFC 2833. In \textsf{Settings} $\to$ \textsf{Audio} $\to$ \textsf{DTMF mode}, force the \texttt{RFC 2833} or \texttt{SIP INFO} method.
  \item Confirm that the \texttt{sip.conf} file has the line \texttt{dtmfmode=rfc2833} for extension 102.
\end{itemize}

\section{Phase wrap-up}

At this point the following is validated:

\begin{itemize}
  \item Asterisk accepts SIP clients and registers them correctly.
  \item The dialplan answers calls and plays audio.
  \item The \texttt{Read()} function captures DTMF into a variable.
  \item The \texttt{UserEvent} is emitted and visible on AMI.
\end{itemize}

That is everything the bridge needs from the Asterisk side. The next phase prepares DOSBox-Staging and Hugo so that the bridge has somewhere to inject the keys.
